\abstract{
Bayesian change-point and segmentation models provide uncertainty-aware piecewise-constant representations of one-dimensional signals, but exact inference is often limited to a narrow set of likelihoods, regular grids, or single-series settings. We present \texttt{BayesBreak}, an application-agnostic framework for fast, exact Bayesian segmentation under weighted exponential-family likelihoods with conjugate priors. \texttt{BayesBreak} reduces each candidate segment to constant-time marginal likelihoods and posterior moments computed from cumulative sufficient statistics, enabling dynamic programming over segmentations to obtain (i) marginal evidences $p(\mathcal{D}\mid k)$, (ii) posterior distributions over boundary locations, and (iii) posterior moments of the latent piecewise-constant signal (a ``Bayesian curve'') together with MAP segmentations (breakpoints and levels). We generalize classic Bayesian piecewise-constant regression to irregularly spaced designs via length-aware segmentation priors, and to pooled inference across replicates and groups, including an EM procedure for latent group assignments. At prediction time we formalize two complementary interfaces: scoring group-membership likelihoods for new data given $(X,Y)$, and evaluating MAP/Bayesian signal summaries given $(X,\widehat{g})$, handling both single pointwise observations and set-valued multivariate units. Finally, we extend the framework to non-conjugate GLMs via local (Laplace/variational) block-evidence approximations and quantify how block-level approximation error propagates through the segmentation DP. The block/DP separation makes BayesBreak extensible: to add a new observation model one only needs a routine for block marginal likelihood (exact or approximate), while the global posterior recursions and prediction machinery remain unchanged. The same interface accommodates both univariate observations and multivariate or multi-channel observations through vector-valued sufficient statistics and segment parameters, enabling principled segmentation and prediction for vector responses under shared or channel-specific boundary structures. We outline efficient implementations and a comprehensive evaluation plan spanning segmentation accuracy, uncertainty calibration, scalability, and group-prediction performance on synthetic benchmarks and real datasets.
}
